% ======================================================================
%  KAVRAM: Context Engineering  (Tier 1)
%  Context Window + Token + Compression burada toplanır
% ======================================================================

\concept{Context · Bağlam}{Context\\Engineering}{Doğru bilgiyi, doğru biçimde, sınırlı pencereye sığacak şekilde yerleştirmek}

% ----------------------------------------------------------------------
\stage{1}{Problem}
% ----------------------------------------------------------------------
Model yalnızca \emph{şu an} context penceresinde olanı bilir. Çok az koyarsan
soruyu cevaplayamaz; çok fazla koyarsan maliyet ve gecikme artar, doğru bilgi
gürültüde kaybolur (\en{lost in the middle}) ve model dağılır.

\keyline{Çoğu ``model kötü'' şikâyeti aslında bir context sorunudur: modele yanlış şeyi, yanlış miktarda, yanlış sırada verdik.}

Context penceresi büyüdükçe sorun çözülmez, yer değiştirir. ``Her şeyi koy''
stratejisi 8K'da da, 1M token'da da aynı tuzağa düşer: alaka, hacimle ölçülmez.

% ----------------------------------------------------------------------
\stage{2}{Neden Var}
% ----------------------------------------------------------------------
Context penceresi, modelin bir çağrı için sahip olduğu \textbf{tüm çalışma
belleğidir}. Model uzun süreli hafızaya, dosya sistemine, canlı veriye sahip
değildir --- ne biliyorsa o an context'tedir. Dolayısıyla:

\begin{itemize}
  \item Pencereye \emph{ne girdiğine} karar vermek, sistemin kalitesini doğrudan belirler.
  \item Pencere sınırlıdır ve her token maliyet/gecikme/dikkat açısından bir \emph{tercihtir}.
  \item ``Context engineering'' bu tercihleri sistemli yapma disiplinidir --- prompt yazmaktan daha geniştir.
\end{itemize}

% ----------------------------------------------------------------------
\stage{3}{Tarihçe}
% ----------------------------------------------------------------------
\begin{timeline}
\yr{2020–22} küçük pencereler (2–4K): bağlam kıt, dikkatle seçilir \tlsep \yr{2023} pencereler büyür (32–128K); ``her şeyi koy'' cazip \tlsep \yr{2023–24} \en{lost in the middle}: model ortayı ihmal eder \tlsep \yr{2024–26} disiplinli montaj: retrieval · özet · sıkıştırma · sıralama
\end{timeline}

% ----------------------------------------------------------------------
\stage{4}{Tanım}
% ----------------------------------------------------------------------
\begin{problembox}
\textbf{Context Engineering:} Her model çağrısı için, bir \emph{token bütçesi}
altında, hangi bilginin --- hangi biçimde ve hangi sırada --- context penceresine
gireceğine karar verme disiplini.

\smallskip
\textbf{Context Window:} Modelin tek çağrıda işleyebileceği toplam token sınırı.
\textbf{Token:} Metnin model tarafından işlenen en küçük birimi; bütçe burada ölçülür.
\end{problembox}

% ----------------------------------------------------------------------
\stage{5}{Zihinsel Model}
% ----------------------------------------------------------------------
\begin{mentalbox}
Context penceresini bir \textbf{çalışma masası} gibi düşün. Yüzeyi sınırlı.
Tüm dosya dolabını masaya boşaltmazsın --- bu işi yavaşlatır ve aradığını
kaybettirir. Tam da bu görev için gereken birkaç belgeyi, en önemlisi elinin
altında olacak şekilde yerleştirirsin.

\smallskip
İyi context engineering, masaya \emph{ne koyduğun} kadar \textbf{neyi dışarıda
bıraktığınla} ilgilidir.
\end{mentalbox}

% ----------------------------------------------------------------------
\stage{6}{Mimari Diyagram}
% ----------------------------------------------------------------------
Token bütçesinin dağılımı (üstte) ve modelin dikkatinin bağlam boyunca nasıl
azaldığı (altta, ``lost in the middle'').

\vspace{4pt}
\begin{center}
\begin{tikzpicture}[font=\sffamily\scriptsize]
  % --- Bütçe çubuğu ---
  \def\H{7mm}
  \fill[accentdk] (0,0) rectangle (18mm,\H);
  \fill[accent!55] (18mm,0) rectangle (54mm,\H);
  \fill[good!55] (54mm,0) rectangle (84mm,\H);
  \fill[softgray!40] (84mm,0) rectangle (100mm,\H);
  \draw[hairline, line width=0.8pt] (0,0) rectangle (118mm,\H);
  \fill[white] (100mm,0) rectangle (118mm,\H);
  \draw[hairline, line width=0.8pt] (100mm,0) rectangle (118mm,\H);

  \node[font=\sffamily\tiny, text=white] at (9mm,\H/2) {System};
  \node[font=\sffamily\tiny, text=white] at (36mm,\H/2) {Geçmiş + getirilen bilgi};
  \node[font=\sffamily\tiny, text=white] at (69mm,\H/2) {User};
  \node[font=\sffamily\tiny, text=ink] at (92mm,\H/2) {çıktı};
  \node[font=\sffamily\tiny, text=softgray] at (109mm,\H/2) {pay};
  \node[anchor=west, font=\sffamily\bfseries\tiny, text=softgray] at (0,\H+3mm)
        {TOKEN BÜTÇESİ \,(sabit toplam)};

  % --- Dikkat eğrisi (U) ---
  \begin{scope}[shift={(0,-22mm)}]
    \draw[hairline, line width=0.8pt, -{Stealth[length=1.6mm]}] (0,0) -- (118mm,0)
      node[right, font=\sffamily\tiny, text=softgray]{};
    \draw[hairline, line width=0.8pt, -{Stealth[length=1.6mm]}] (0,0) -- (0,14mm)
      node[above, font=\sffamily\tiny, text=softgray]{dikkat};
    \draw[accent, line width=1.1pt]
      (2mm,12mm) .. controls (30mm,1mm) and (88mm,1mm) .. (116mm,12mm);
    \node[font=\sffamily\tiny, text=ink] at (8mm,13.5mm) {baş};
    \node[font=\sffamily\tiny, text=warn] at (59mm,-2.5mm) {orta \,(kaybolur)};
    \node[font=\sffamily\tiny, text=ink] at (110mm,13.5mm) {son};
  \end{scope}
\end{tikzpicture}
\end{center}

\keyline{Sonuç: kritik bilgiyi bağlamın başına veya sonuna koy; ortaya gömme. Ve her token'ı hak ettiği için harca.}

% ----------------------------------------------------------------------
\stage{7}{Karar Çerçevesi}
% ----------------------------------------------------------------------
Bir bilgi parçası context'e girmeli mi, girmeliyse hangi biçimde? Token bütçesi
karşısında her parça için aynı soruları sor.

\vspace{2pt}
\begin{center}
\begin{tikzpicture}[
   font=\sffamily\scriptsize, node distance=5mm and 8mm,
   q/.style={draw=accent, line width=0.8pt, fill=wash, diamond, aspect=2.4,
             align=center, inner sep=1pt, minimum width=30mm},
   leaf/.style={draw=accentdk, line width=0.8pt, rounded corners=2pt,
                fill=accent!12, align=center, inner sep=4pt, minimum width=22mm},
   yes/.style={good, font=\sffamily\bfseries\scriptsize},
   no/.style={warn, font=\sffamily\bfseries\scriptsize},
   ed/.style={-{Stealth[length=2mm]}, softgray, line width=0.7pt}]

  % dikey omurga: sorular aşağı iner, yapraklar sola/aşağı dallanır
  \node[q] (q1) {Bu görev için\\gerekli mi?};
  \node[leaf, left=14mm of q1] (drop) {\textbf{Dışarıda bırak}};
  \node[q, below=11mm of q1] (q2) {Bütçeye sığıyor\\mu?};
  \node[leaf, left=14mm of q2] (incl) {\textbf{Olduğu gibi koy}};
  \node[q, below=11mm of q2] (q3) {Tamamı mı\\gerekli?};
  \node[leaf, below left=12mm and 4mm of q3, minimum width=20mm] (sum) {\textbf{Özetle}\\sıkıştır};
  \node[leaf, below right=12mm and 4mm of q3, minimum width=20mm] (retr) {\textbf{Retrieve}\\en iyi k};

  \draw[ed] (q1) -- node[no,above,inner sep=1pt]{Hayır} (drop);
  \draw[ed] (q1) -- node[yes,right]{Evet} (q2);
  \draw[ed] (q2) -- node[yes,above,inner sep=1pt]{Evet} (incl);
  \draw[ed] (q2) -- node[no,right]{Hayır} (q3);
  \draw[ed] (q3) -- node[yes,left]{Evet} (sum);
  \draw[ed] (q3) -- node[no,right]{Hayır} (retr);
\end{tikzpicture}
\end{center}

\keyline{Sıra: önce alaka (gerekli mi?), sonra biçim (olduğu gibi / özet / retrieval). Hacim en son gelir.}

% ----------------------------------------------------------------------
\stage{8}{Üretim Senaryosu}
% ----------------------------------------------------------------------
\begin{prodbox}[Üretim Senaryosu · Uzayan Sohbet]
\textbf{Bağlam.} Bir asistan, uzun sohbetlerde tüm geçmişi her çağrıda context'e
ekliyor. 30. mesajdan sonra cevaplar yavaşlıyor, pahalılaşıyor ve ilk kuralları
``unutuyor'' (çünkü kurallar bağlamın ortasında kalıyor).

\smallskip
\textbf{Çözüm.} Eski turlar tek bir \emph{rolling summary}'ye sıkıştırılır;
yalnızca son birkaç tur ham tutulur; gerektiğinde geçmişten ilgili parçalar
retrieve edilir. System kuralları her zaman \emph{başta} kalır.

\smallskip
\textbf{Sonuç.} Token kullanımı sabitlenir, gecikme düşer, tutarlılık geri gelir ---
pencere büyütülmeden. Sorun hacimde değil, montajdaymış.
\end{prodbox}

% ----------------------------------------------------------------------
\stage{9}{Anti-pattern}
% ----------------------------------------------------------------------
\begin{antibox}[Anti-pattern · ``Infinite Context / Her Şeyi Koy'']
``Pencere büyük, hepsini koyalım, model halleder.'' Tüm dokümanları, tüm geçmişi,
her ihtimali context'e doldurmak. Sonuç: dikkat seyrelir, kritik bilgi ortada
kaybolur, maliyet ve gecikme patlar, model alakasız ayrıntılarla dağılır.

\smallskip
\textbf{Belirti:} ``Daha çok bağlam verdik ama cevaplar kötüleşti / yavaşladı.''\\
\textbf{Kök neden:} Hacmi alakayla karıştırmak; bütçeyi tercih olarak görmemek.\\
\textbf{Düzeltme:} Alakaya göre seç, özetle/retrieve et, kritik bilgiyi uçlara koy, kullanılmayanı çıkar.
\end{antibox}

% ----------------------------------------------------------------------
\stage{10}{Alternatifler \& Ödünleşimler}
% ----------------------------------------------------------------------
\vspace{2pt}
\begin{tabularx}{\linewidth}{@{}>{\raggedright\arraybackslash}p{30mm} L L@{}}
\arrayrulecolor{hairline}\toprule
{\sffamily\bfseries\small Strateji} &
{\sffamily\bfseries\small Ne zaman} &
{\sffamily\bfseries\small Ödünleşim} \\
\midrule
\textbf{Daha büyük pencere} & Bilgi gerçekten büyük ve hepsi gerekli &
Pahalı, yavaş, lost-in-the-middle \\
\addlinespace[3pt]
\textbf{Retrieval (RAG)} & Bilgi çok büyük, sadece ilgili parça gerek &
Retrieval altyapısı; isabet riski \\
\addlinespace[3pt]
\textbf{Özetleme} & Geçmiş/uzun metin sıkıştırılabilir &
Detay kaybı; özet hatası birikir \\
\addlinespace[3pt]
\textbf{Sıkıştırma/budama} & Tekrar/gürültü çok &
Yanlış budama bilgi kaybettirir \\
\bottomrule
\end{tabularx}

\vspace{4pt}
Gerçek sistemler bunları \emph{birlikte} kullanır: sabit system + özetlenmiş
geçmiş + retrieve edilmiş bilgi + ham user isteği.

% ----------------------------------------------------------------------
\stage{11}{İlişkili Kavramlar}
% ----------------------------------------------------------------------
\begin{itemize}
  \item \term{Context Window} --- montajın sığması gereken sert sınır.
  \item \term{Token} --- bütçenin para birimi.
  \item \term{RAG} --- pencereye yalnızca ilgili bilgiyi getirme yöntemi.
  \item \term{Summarization} --- geçmişi/uzun metni sıkıştırma.
  \item \term{Prompt} --- montajın görünen yüzü; context onu besler.
\end{itemize}

% ----------------------------------------------------------------------
\stage{12}{Özet}
% ----------------------------------------------------------------------
\begin{itemize}
  \item Context penceresi modelin \textbf{tüm çalışma belleğidir}; ne koyarsan onu bilir.
  \item Pencere büyütmek sorunu çözmez, yer değiştirir: alaka hacimle ölçülmez.
  \item Token bir bütçedir; her token bir tercihtir.
  \item Kritik bilgiyi uçlara koy --- ortası ``kaybolur''.
  \item Seç → biçimlendir (özet/retrieval) → sırala. Hacim en son düşünülür.
\end{itemize}

% ----------------------------------------------------------------------
\stage{13}{Mühendislik İlkesi}
% ----------------------------------------------------------------------
\begin{principlebox}
\textbf{İlke \#3 — Context bir bütçedir; her token bir tercihtir.}\\[4pt]
\normalsize Modele her şeyi verme; \emph{tam da gerekeni} ver. Kaliteyi
belirleyen, modelin ne kadar gördüğü değil, \emph{ne gördüğüdür}.
\end{principlebox}
